\chapter{Conclusion et perspectives}
\label{ch:conclusion}

\section{Bilan}

Le projet \stockflow{} démontre qu'une application métier de gestion de stock peut être livrée avec une chaîne DevOps complète : développement local Docker Compose (avec Adminer), validation Kubernetes sur minikube, observabilité Prometheus/Grafana, visualisation du cluster et CI/CD GitHub Actions — le tout \textbf{sans serveur de production obligatoire}.

\section{Objectifs atteints}

\begin{itemize}
  \item Fonctionnalités métier : catalogue, mouvements, alertes, tableau de bord ;
  \item Développement local via \texttt{make dev} (Adminer pour PostgreSQL) ;
  \item Validation K8s via \texttt{make cluster} (minikube, Helm, Grafana, dashboard) ;
  \item CI avec six jobs dont déploiement éphémère minikube ;
  \item Infrastructure as Code (Terraform validate en CI ; apply VPS optionnel) ;
  \item Déploiement production documenté (\texttt{deploy.yml} manuel, \texttt{docs/vps-setup.md}).
\end{itemize}

\section{Limites connues}

\begin{itemize}
  \item Cluster mono-nœud (minikube ou k3s) ;
  \item State Terraform local ;
  \item Tests CI encore majoritairement SQLite ;
  \item Monitoring complet (kube-prometheus) non exécuté dans le job CI minikube (RAM) — preuve en démo locale ;
  \item Identifiants de démo à renforcer en production réelle.
\end{itemize}

\section{Perspectives}

\begin{enumerate}
  \item Backend Terraform distant (S3 + DynamoDB) ;
  \item Haute disponibilité multi-nœuds ;
  \item GitOps (Flux / Argo CD) ;
  \item Tests e2e post-déploiement ;
  \item Notifications sur alertes stock ;
  \item Déploiement VPS stabilisé si ressources matérielles suffisantes.
\end{enumerate}

\section{Mot de fin}

\stockflow{} illustre l'unification code applicatif et infrastructure dans un même dépôt, avec deux chemins de validation complémentaires — Docker Compose pour le développement, minikube pour l'orchestration — et une CI qui garantit la reproductibilité du déploiement Helm.
